\section{Caso de un preesquema base normal}
\label{V.8}
\marginpar{143}%
\begin{proposition}
\label{V.8.1}
Sea~$S$ el espectro de un cuerpo~$k$, y sea~$\Omega$ una extensión
algebraicamente cerrada de~$k$ que define un punto geométrico~$a$ de~$S$
con valores en~$\Omega$. Sea~$\bar{k}$ la clausura separable de~$k$
en~$\Omega$. Entonces existe un isomorfismo canónico de~$\pi_1(S,a)$
sobre el grupo de Galois topológico de~$\bar{k}/k$.
\end{proposition}

Sea~$k'$ la clausura algebraica de~$k$ en~$\Omega$; corresponde, pues, a un
punto geométrico~$b$ de~$S$, con valores en~$k'$. El homomorfismo natural
de funtores~$F_b\to F_a$ es evidentemente un isomorfismo, pues todo
$k$-homomorfismo de una extensión finita separable de~$k$ en~$\Omega$
toma necesariamente sus valores en~$\bar{k}$ y, con mayor razón, en~$k'$.
Por otra parte, el grupo~$\pi'$ de los $k$-automorfismos de~$k'$ actúa de
manera evidente sobre~$F_b$, de donde se obtiene un homomorfismo
$\pi'\to\Aut(F_b)\isomto\Aut(F_a)=\pi_1(S;a)$.
Por otra parte, es bien sabido que el homomorfismo natural de~$\pi'$ en el
grupo~$\pi$ de los automorfismos de~$\bar{k}/k$ es un isomorfismo. Se
obtiene así un homomorfismo canónico~$\pi\to\pi_1(S;a)$; queda por demostrar
que es un isomorfismo. En efecto, este homomorfismo es inyectivo, pues un
elemento del núcleo es un automorfismo de~$\bar{k}/k$ que induce la identidad
en toda subextensión finita separable y, por tanto, es trivial. Este
homomorfismo es sobreyectivo, pues si~$X$ es un recubrimiento étale
\emph{conexo} de~$S$, definido, por tanto, por una \emph{extensión} finita
separable~$L/k$, entonces~$\pi$ actúa transitivamente sobre el conjunto de
los $k$-homomorfismos de~$L$ en~$k'$,
\ifthenelse{\boolean{orig}}{como bien sabido}{como es bien sabido}.

\begin{proposition}
\label{V.8.2}
Sean~$S$ un preesquema conexo, localmente noetheriano y normal,
$K=\kres(s)$ su cuerpo de funciones $=$ el cuerpo residual en su punto
genérico~$s$, y~$\Omega$ una extensión algebraicamente cerrada de~$K$ que
define un punto geométrico~$a'$ de~$S'=\Spec(K)$ y un punto geométrico~$a$
de~$S$. Entonces el homomorfismo
$\pi_1(S';a')\to\pi_1(S;a)$ es sobreyectivo. Cuando se identifica el primer
grupo con el grupo de Galois de la clausura separable~$\bar{K}$ de~$K$
en~$\Omega$ (\cf~\Ref{V.8.1}), el núcleo del homomorfismo anterior
corresponde, por la teoría de Galois, a la subextensión de~$\bar{K}/K$ que
es el compuesto de las extensiones finitas de~$K$ en~$\Omega$ no ramificadas
sobre~$S$.
\end{proposition}

La primera afirmación significa que la imagen inversa sobre~$S'$ de un
recubrimiento étale conexo~$X$ de~$S$ es conexa, \ie que~$X$ es íntegro;
esto no es sino~\Ref{I.10.1}. El núcleo del homomorfismo anterior se
interpreta entonces como el conjunto de los automorfismos de~$\bar{K}/K$
que inducen la identidad sobre los conjuntos~$F_a(X)$,
\marginpar{144}%
donde puede suponerse conexo el recubrimiento étale~$X$ de~$S$. Pero esto
significa que tal automorfismo induce la identidad en las subextensiones
finitas de~$\bar{K}/K$ no ramificadas sobre~$S$, lo que demuestra la última
afirmación.

\begin{remarquestar}
Gracias a esta interpretación del grupo fundamental del preesquema normal~$S$
en términos de la
\ifthenelse{\boolean{orig}}{teoría de los Galois usual}{teoría de Galois usual},
la definición se conocía desde hacía mucho tiempo en este caso.
\end{remarquestar}
